% !TeX root = ../../Book4.tex
% 纲要标题：## 附录 A　Stein 空间与解析层
% 纲要位置：几何学计划纲要.md，第 4598 行。
% 正文状态：未撰写；本文件目前是框架。

% 纲要内容开始
% 
% > 局部多复变之后的独立整体路线。正文需要的凝聚性不留作附录黑箱。
% 
% * **Stein 条件与模型**
%   * 全纯凸性和可分性
%   * 严格 psh 穷竭
%   * 仿射簇、域及非 Stein 模型
% 
% * **Cartan 方法**
%   * 定理 A、B 的准确对象
%   * Cousin 问题与上同调消失
%   * Oka 原理的所选陈述
% 
% * **∂̄ 与整体嵌入**
%   * 带权估计和穷竭
%   * 函数分离及嵌入问题
%   * 奇异 Stein 空间的额外工具
% 
% ---
% 
% 纲要内容结束

\chapter{Stein 空间与解析层}
\label{g4:app:A}
\BookChapterNumber{g4:app:A}

待撰写

\section{Stein 条件与模型}
\label{g4:sec:A01}
待撰写

\subsection{全纯凸性和可分性}
\label{g4:subsec:A01:01}
待撰写

\subsection{严格 psh 穷竭}
\label{g4:subsec:A01:02}
待撰写

\subsection{仿射簇、域及非 Stein 模型}
\label{g4:subsec:A01:03}
待撰写

\section{Cartan 方法}
\label{g4:sec:A02}
待撰写

\subsection{定理 A、B 的准确对象}
\label{g4:subsec:A02:01}
待撰写

\subsection{Cousin 问题与上同调消失}
\label{g4:subsec:A02:02}
待撰写

\subsection{Oka 原理的所选陈述}
\label{g4:subsec:A02:03}
待撰写

\section{\texorpdfstring{\(\bar{\partial}\) 与整体嵌入}{\={∂} 与整体嵌入}}
\label{g4:sec:A03}
待撰写

\subsection{带权估计和穷竭}
\label{g4:subsec:A03:01}
待撰写

\subsection{函数分离及嵌入问题}
\label{g4:subsec:A03:02}
待撰写

\subsection{奇异 Stein 空间的额外工具}
\label{g4:subsec:A03:03}
待撰写
